% =============================================================================
% Daftar Pustaka
%
% PPTA IPB 2026: Harvard nama-tahun dengan gaya CSE edisi ke-9:
%   - In-text: (Author Year) atau (Author1 dan Author2 Year) atau
%              (Author1 et al. Year) untuk >2 penulis
%   - Daftar pustaka urut alfabet, hanging indent 1 cm, periode sebagai
%     pemisah field, tahun tanpa kurung.
%   - Maksimum lima penulis, lalu "et al." untuk sumber dengan >5 penulis.
%   - DOI dicantumkan bila tersedia; sumber daring memuat tanggal akses dan URL.
%
% Perintah yang bisa dipakai di teks (sintaks natbib kompatibel):
%   \citep{key}                  -> (Author Year)
%   \citep[5]{key}               -> (Author Year:5)
%   \citet{key}                  -> Author (Year)
%   \citep{key1,key2}            -> (Author1 Year1; Author2 Year2)
%   \citeauthor{key}             -> Author
%   \citeyear{key}               -> Year
%
% Backend: biber. Style didefinisikan di src/config/ipb-thesis.sty.
% Daftar pustaka di-deklarasikan di src/main.tex melalui \addbibresource.
% =============================================================================

\chapter*{DAFTAR PUSTAKA}
\addcontentsline{toc}{chapter}{DAFTAR PUSTAKA}

% Cetak daftar pustaka. Heading none agar tidak ada judul tambahan
% (judul "DAFTAR PUSTAKA" sudah ditangani oleh \chapter* di atas).
% Baris berikut menampilkan lima contoh dari DOTX. Hapus setelah contoh
% diganti; pada naskah akhir hanya sumber yang benar-benar diacu yang dicetak.
\nocite{Arifin2023,Aulia2024,Brett2025,Pezzino2026,Rajam2022}
\printbibliography[heading=none]
